\scanpage{67}
\mspage{33}
``drilling holes'' must provide us with an outer group $\Pi_{0,3}(I)$, \inserted{of type $(0,3+\operatorname{card} I)$,} on which $\Gamma$ acts (it is perhaps better to use the yoga of rigid groupoids introduced elsewhere — so one \unsure{can|will be able to|for} \unsure{or|in} will also have to plug some of the holes $0,1,\infty$ — \unsure{one|in} therefore finds the groupoid equivalent of the \deleted{``\unsure{arithmetic|arithmetic} sph\unsure{e|e}re} projective line $\mathbb P^1_Q$, one calls it the arithmetic sphere — which corresponds to an infinite \inserted{outer} group with loops on which $\Gamma$ acts — in fact, it is none other than $E_{K_1}$, where $\vert$
\[
\tag{6}K_1=Q(\deleted{\unsure{X|x}_{\unsure{1|4}}}\,T_1)
\]
is the \inserted{typical} purely transcendental extension of degree $1$ of $Q$.

Starting from (6), one constructs from $\overline m$ the groupoid equivalent of $\deleted{\unsure{X|P}}\,U_{0,3}$ \unsure{$\longleftarrow$|\deleted{(\unsure{qu}}} and one \unsure{goes back up|take|will take} as
\margin{66}
